\documentclass[11pt]{article}
\usepackage[margin=1in]{geometry}
\usepackage[T1]{fontenc}
\usepackage{lmodern,microtype,hyperref,enumitem,amsmath,amssymb,xcolor,booktabs}
\usepackage[numbers,sort&compress]{natbib}
\hypersetup{colorlinks=true,urlcolor=blue,citecolor=blue,linkcolor=blue}
\title{Q05: When Does Extra Reasoning Actually Improve an Answer?}
\author{DLS 2026 --- module author: Kukla}
\date{October 2026}
\usepackage{url}
\begin{document}
\maketitle

\section*{Learning objectives}
\begin{itemize}[leftmargin=1.4em]
\item Explain what ``extra reasoning'' (chain-of-thought, test-time compute) actually changes at inference.
\item State the conditions under which longer reasoning helps, hurts, or is neutral.
\item Reason quantitatively about the accuracy-vs-cost trade-off of test-time compute.
\end{itemize}

\begin{quote}\small\textbf{Companion video:}\begin{itemize}\setlength\itemsep{0pt}\item \emph{OpenAI o1's New Paradigm: Test-Time Compute Explained} --- bycloud (15:43). \url{https://www.youtube.com/watch?v=pi7LF-OpO6k}\end{itemize}\end{quote}

\section{What ``extra reasoning'' means}
Student response \#05 asks when extra reasoning actually improves an answer and what it costs.
``Extra reasoning'' here means spending more inference-time computation before committing to a final
answer: emitting intermediate steps (\emph{chain-of-thought}, CoT \citep{wei2022cot}), sampling many
candidate solutions and aggregating (\emph{self-consistency} \citep{wang2023selfconsistency}), or using
a trained reasoning model that produces long deliberate traces (o1-style / R1-style
\citep{deepseekr1}). All of these trade tokens---hence latency and money---for a chance at higher
accuracy. The central question is \emph{when that trade is worth it}.

\section{When extra reasoning helps}
The gains are \textbf{task-dependent}, and this is the key teaching point.
\begin{itemize}[leftmargin=1.4em]
\item \textbf{Multi-step problems} (math, code, logic, planning) benefit most: CoT lets the model
externalize intermediate state it would otherwise have to compute in a single forward pass
\citep{wei2022cot}. Empirically the benefit grows with the number of reasoning steps the problem
genuinely requires.
\item \textbf{Problems with a verifiable answer} benefit from sampling-and-selecting: self-consistency
and best-of-$n$ with a verifier convert extra samples into accuracy, because a correct answer only has
to be \emph{found} once and then recognized \citep{wang2023selfconsistency,cobbe2021verifiers}.
\item \textbf{Test-time compute can substitute for parameters}: on some benchmarks, letting a smaller
model think longer matches a larger model answering immediately \citep{snell2024scaling}.
\end{itemize}

\section{When it does not help---or hurts}
\begin{itemize}[leftmargin=1.4em]
\item \textbf{Simple or retrieval-like tasks}: for single-step factual lookups, CoT adds tokens and
latency with little or no accuracy gain, and can even \emph{lower} accuracy by talking the model out of
a correct first instinct.
\item \textbf{Overthinking}: forcing long traces on easy problems wastes compute and raises the chance
of a reasoning error compounding. Recent work documents an inverted-U: accuracy rises then falls as
reasoning length grows past what the task needs.
\item \textbf{Faithfulness gap}: the written reasoning is not guaranteed to be the computation that
produced the answer \citep{turpin2023faithfulness}. A plausible-looking trace can accompany a wrong
answer, which matters when we use the trace to \emph{trust} the output.
\end{itemize}

\section{The trade-off, quantitatively}
Let $p_1$ be single-sample accuracy and suppose we draw $n$ independent samples and select with a
verifier of recall $r$. Coverage (probability at least one sample is correct) is
$1-(1-p_1)^n$, so best-of-$n$ accuracy is upper-bounded by that and, with a perfect verifier, equals it.
Cost scales roughly linearly in $n$ (and in trace length for CoT). The practical consequence: returns
\emph{diminish}---each extra sample buys less marginal accuracy---so there is an economically optimal
compute budget per query, and it depends on the value of a correct answer versus the price of tokens.
The companion code (\texttt{code/best\_of\_n.py}) simulates this coverage curve.

\section{A decision rule for practitioners}
Spend extra reasoning when: (1) the task is genuinely multi-step, (2) errors are costly, and (3) you can
\emph{verify} or at least cross-check candidate answers. Keep it short when the task is single-step,
latency-sensitive, or unverifiable. Calibrate the budget per task---do not apply one reasoning length
globally.

\section{Summary}
Extra reasoning trades inference compute for accuracy. It pays on multi-step, verifiable problems and
is wasteful or harmful on simple ones. Returns diminish with sample count and trace length, the written
trace may be unfaithful, and the right budget is task-specific---so the honest answer to ``does
reasoning help?'' is ``it depends, measurably, on the task.''

\bibliographystyle{plainnat}
\bibliography{bibliography}
\end{document}
